%%%%%%%%%%  scan idx 634  |  folio 621  %%%%%%%%%%
\SGAPrintPage{621}%
with \(G\) a \((b_1,\ldots,b_r)\)-sheaf and \(\operatorname{dim{.}supp}(G)=s-1\). Then
\[
h^0(F(n))-h^0(F(n-1))\le h^0(G(n)),
\]
but
\[
h^0(G(n))\le\sum_{i=0}^{s-1}b_{i+1}\binom{n+i}{i}
\]
by induction on \(s\) and \(h^0(F(-1))\le b_0\); hence (i) by induction on \(n\ge -1\). Finally, by 1.7 and induction on \(s\), \(a_s\le b_s\) and \(G\) is a \((b_1,\ldots,b_{s-1},a_s)\)-sheaf; hence (ii).

\sgaStableTarget{sga6.e13.def.1-9}\textbf{Definition 1.9.} \emph{The following polynomials with rational coefficients \(\ge0\), defined by induction, are called the \emph{\((b)\)-polynomials}}~:
\[
\begin{cases}
P_{-1}=0\\
P_r(X_0,\ldots,X_r)=P_{r-1}(X_1,\ldots,X_r)
+\displaystyle\sum_{i=0}^{r}X_i
\binom{P_{r-1}(X_1,\ldots,X_r)-1+i}{i}.
\end{cases}
\]

\sgaStableTarget{sga6.e13.rem.1-10}\textbf{Remark 1.10.} For the \((b)\)-polynomials, one sees at once by induction on \(r\ge t\) that
\[
P_r(X_0,\ldots,X_t,0,\ldots,0)=P_t(X_0,\ldots,X_t).
\]

\sgaStableTarget{sga6.e13.thm.1-11}\textbf{Theorem 1.11.} \emph{(The main theorem on \((b)\)-sheaves). Let \(F\) be a \((b)\)-sheaf on \(X\) and \(\chi(F(n))=\sum_{i=0}^{r}a_i\binom{n+i}{i}\) the Hilbert polynomial of \(F\); let \((c)=(c_0,\ldots,c_r)\) be a sequence of integers \(\ge (b)-(a)\) and \(m=P_r(c)\) the integer given by the \(r\)-th \((b)\)-polynomial (1.9). Then \(m\ge0\) and \(F\) is \(m\)-regular. In particular, if \(s=\operatorname{dim{.}supp}(F)\), then \(F\) is \(P_{s-1}(c_0,\ldots,c_{s-1})\)-regular.}

%%%%%%%%%%  scan idx 635  |  folio 622  %%%%%%%%%%
\SGAPrintPage{622}%
Indeed, we argue by induction on \(r\). When \(r=0\), this is clear since \(m\ge0\) and \(F\) is \(n\)-regular for every \(n\). When \(r\ge1\), we have an exact sequence $0\longrightarrow F(-1)\longrightarrow F\longrightarrow G\longrightarrow0$ where \(G\) is a \((b_1,\ldots,b_r)\)-sheaf. By 1.7 and the induction hypothesis, we see that if \(n=P_{r-1}(c_1,\ldots,c_r)\), then \(n\ge0\) and \(G\) is \(n\)-regular. Thus, by 1.4, \(F\) is \([n+h^1(F(n-1))]\)-regular and \(h^q(F(n-1))=0\) for \(q\ge2\). Therefore, \(h^1(F(n-1))=h^0(F(n-1))-\chi(F(n-1))\), which is \(\le \sum_{i=0}^{r}(b_i-a_i)\binom{n-1+i}{i}\) by 1.8 (i), since \(b_i\ge0\). Finally, \(F\) is \(\left[n+\sum_{i=0}^{r}c_i\binom{n-1+i}{i}\right]\)-regular by 1.3 (i); this gives the first assertion. The second follows from it by 1.8 (ii) and 1.10.

\sgaStableTarget{sga6.e13.para.1-12}\textbf{1.12.} Let \(S\) be a noetherian scheme and \(X\) an \(S\)-scheme of finite type. Let \(\mathcal F\) be a \uline{family of classes of coherent sheaves on the fibers of} \(X/S\); that is, for every point \(s\in S\) and every extension \(K\) of \(k(s)\), coherent sheaves \(F_K\) on the algebraic scheme \(X_K\) are given, and we say that an \(F_K\) and an \(F'_{K'}\) belong to the same class if there exist \(k(s)\)-homomorphisms from \(K,K'\) into a common extension \(K''\) of \(k(s)\) such that \(F_{K''}(=F_K\otimes_K K'')\) and \(F'_{K''}\) are isomorphic on \(X_{K''}\). Then the family \(\mathcal F\) is said to be \uline{bounded by the coherent sheaf} \(F\) \uline{on} \(X_T=X\times_S T\), where \(T\) is of finite type over \(S\), if \(\mathcal F\) is contained in the family of classes of the sheaves \(F_{k(t)}\) with \(t\in T\); we say that \(\mathcal F\) is \uline{bounded} if there exist \(T,F\) as above.

Suppose in addition that \(X/S\) is projective and endowed with a (relatively) ample invertible sheaf \({O}_X(1)\). Then \(\mathcal F\) is said to be a \((b)\)-\uline{family}, for a sequence of integers \((b)=(b_0,\ldots,b_r)\), if every class in \(\mathcal F\) is represented by an \(F_K\) with \(K\) algebraically closed, which is a \((b)\)-sheaf.

%%%%%%%%%%  scan idx 636  |  folio 623  %%%%%%%%%%
\SGAPrintPage{623}%
\sgaStableTarget{sga6.e13.thm.1-13}\textbf{Theorem 1.13.}
\emph{Let \(S\) be a noetherian scheme and \(X\) a projective \(S\)-scheme endowed with an ample sheaf \({O}_X(1)\) whose induced sheaves \({O}_{X_s}(1)\), with \(s\in S\), are generated by their \(H^0({O}_{X_s}(1))\); let \(\mathcal F\) be a family of classes of coherent sheaves on the fibers of \(X/S\). Then the following conditions are equivalent}:
\begin{enumerate}[label=(\roman*)]
\item \emph{\(\mathcal F\) is bounded. Moreover, if all the \(F_K\in\mathcal F\) are locally free of rank \(p\), then one may suppose that \(\mathcal F\) is bounded by a sheaf \(F\) on an \(X_T\), with \(F\) locally free of rank \(p\).}
\item \emph{The set of Hilbert polynomials \(\chi(F_K(n))\) of the \(F_K\in\mathcal F\) is finite, and there exists a sequence of integers \((b)\) such that \(\mathcal F\) is a \((b)\)-family.}
\item \emph{The set of Hilbert polynomials \(\chi(F_K(n))\) of the \(F_K\in\mathcal F\) is finite, and there exists an integer \(m\) such that all the \(F_K\in\mathcal F\) are \(m\)-regular.}
\item \emph{The set of Hilbert polynomials \(\chi(F_K(n))\) of the \(F_K\in\mathcal F\) is finite, and \(\mathcal F\) is contained in the family of classes of quotients of sheaves of the form \(E_K\), where \(E\) is a coherent sheaf on an \(X_T\). Moreover, one may suppose \(T=S\) and \(E\) of the form \({O}_X(-m)^{\oplus M}\).}
\item \emph{\(\mathcal F\) is contained in the family of classes of cokernels of homomorphisms of the form \(E'_K\to E_K\), where \(E'\) and \(E\) are coherent sheaves on an \(X_T\). Moreover, one may suppose \(T=S\) and \(E'\) and \(E\) of the forms \({O}_X(-m')^{\oplus M'}\) and \({O}_X(-m)^{\oplus M}\).}
\end{enumerate}
Indeed, suppose that \(\mathcal F\) is bounded by a sheaf \(F\) on an \(X_T\); after decomposing \(T\) into pieces, using the generic flatness theorem (EGA IV 6.9.1), one may suppose \(F\) flat over \(T\). Then,
%%%%%%%%%%  scan idx 637  |  folio 624  %%%%%%%%%%
\SGAPrintPage{624}%
the number of polynomials \(\chi(F(n))\) is at most equal to the number of connected components of \(T\) (EGA III 7.9.5), and it is a small lemma that if, for a \(t\in T\), \(F_t\) is locally free of rank \(p\), then so is \(F\) above an open neighborhood of \(t\) in \(T\) (see EGA IV 11.3.10, where one takes \(X=Y\), \(f=\mathrm{id}\)). Now it follows from the hypotheses and (EGA III 7.7.6) that, after decomposing \(T\) more finely, one can find a sequence \(\sigma\) of sections \(\sigma_1,\ldots,\sigma_r\in H^0({O}_X(1))\) which is \(F\)-regular. Thus (ii) follows from the upper semi-continuity property of the function \(t\longmapsto h^0(F_{t,i}(-1))\) where, for \(t\in T\) and \(0\le i\le r\), \(F_{t,i}\) denotes the restriction of \(F_t\) to the zero scheme of \(\sigma_1,\ldots,\sigma_i\) (EGA III 7.7.5).

The implication (ii) \(\Rightarrow\) (iii) follows immediately from 1.11. The implication (iii) \(\Rightarrow\) (iv) follows from 1.3 (iii) by taking \(M=\max(\chi(F_K(m)))\) and \(E={O}_X(-m)^{\oplus M}\).

Suppose that \(\mathcal F\) satisfies (iv); there exist exact sequences
\[
0\longrightarrow F'_K\longrightarrow E_K\longrightarrow F_K\longrightarrow0
\]
and the set of polynomials \(\chi(F_K(n))\) is finite. By hypothesis, the family of classes of the \(E_K\) is bounded; thus, by the implication (i) \(\Rightarrow\) (ii) already proved, the set of polynomials \(\chi(E_K(n))\) is finite and there exists a sequence of integers \((b)\) such that all the \(E_K\) (with \(K\) algebraically closed) are \((b)\)-sheaves. Consequently, the set of polynomials \(\chi(F'_K(n))\) is finite and, by 1.6, all the \(F'_K\) are \((b)\)-sheaves. Thus, applying to the family of classes of the \(F'_K\) the implication (ii) \(\Rightarrow\) (iv) already proved, we obtain (v).

Finally, suppose that \(\mathcal F\) satisfies (v); let us prove that \(\mathcal F\) is bounded. One may suppose \(E\) (resp. \(E'\)) of the form \({O}_X(-m)^{\oplus M}\). Indeed,
%%%%%%%%%%  scan idx 638  |  folio 625  %%%%%%%%%%
\SGAPrintPage{625}%
by (i) \(\Rightarrow\) (iv), take a surjection \(L={O}_{X_T}(-m)^{\oplus M}\longrightarrow E\). After decomposing \(T\) into pieces so that \(E\) becomes flat over \(T\), the formation of the exact sequence \(0\longrightarrow I\xrightarrow{u}L\longrightarrow E\longrightarrow0\) commutes with passage to the fibers. By (i) \(\Rightarrow\) (iii) and (i) \(\Rightarrow\) (iv), take \(m_1\) large enough that all the \(\operatorname{Ext}^1({O}_{X_K}(-m_1),I_K)=H^1(I_K(m_1))\) vanish and that there exist surjections \(L_1={O}_{X_T}(-m_1)^{\oplus M_1}\longrightarrow E'\) and \(\alpha:L_2={O}_{X_T}(-m_1)^{\oplus M_2}\longrightarrow I\). Then the maps \(\operatorname{Hom}(L_{1,K},L_K)\longrightarrow \operatorname{Hom}(L_{1,K},E_K)\) are surjective. Let \(\beta:E'_K\to E_K\) be a homomorphism and let \(\gamma\) be the composition \(L_{1,K}\to E'_K\to E_K\); then \(\gamma\) comes from a homomorphism \(\delta:L_{1,K}\to L_K\), and $(\delta,u_K\circ\alpha):L_{1,K}\oplus L_{2,K}\longrightarrow L_K$ and \(\beta:E'_K\to E_K\) have the same cokernel. Consequently, \(\mathcal F\) is contained in the family of classes of cokernels of homomorphisms of the form \({O}_{X_K}(-m_1)^{\oplus(M_1+M_2)}\longrightarrow {O}_{X_K}(-m)^{\oplus M}\).

Again after decomposing \(T\) into pieces, one may suppose in addition that \(X_T\) (hence \(\underline{\Hom}_{{O}_{X_T}}(E',E)\)) and all the \(R^q(f_T)_*\underline{\Hom}_{{O}_{X_T}}(E',E)\) (where \(f_T:X_T\to T\) is the structural morphism) are flat over \(T\), so that the formation of \(G=(f_T)_*\underline{\Hom}_{{O}_{X_T}}(E',E)\) commutes with arbitrary base change \(T'\to T\) (EGA III 6.9.9.2). Let $R=\mathbb V(G^*)$; then \(X_R\) canonically carries a ``universal'' exact sequence
\[
E'_R\longrightarrow E_R\longrightarrow F\longrightarrow0
\]
and \(F\) bounds \(\mathcal F\).

\sgaStableTarget{sga6.e13.sec.2}\subsection*{2. Several technical lemmas}
\sgaStableTarget{sga6.e13.para.2-1}\textbf{2.1.} Let \(k\) be a field and \(X\) a proper \(k\)-scheme. Let \(L_1,\ldots,L_s\) be invertible sheaves on \(X\), \(c_1(L_i)\) the Chern class of \(L_i\) and \(Y\) a closed subscheme of \(X\) of dimension \(\le s\). Recall that we denote by
